\documentclass[11pt]{article}
\usepackage[margin=1.1in]{geometry}
\usepackage{amsmath,amssymb,amsthm}
\usepackage[colorlinks=true,linkcolor=blue,citecolor=blue]{hyperref}

\newtheorem{theorem}{Theorem}[section]
\newtheorem{lemma}[theorem]{Lemma}
\newtheorem{proposition}[theorem]{Proposition}
\newtheorem{corollary}[theorem]{Corollary}
\newtheorem{remark}[theorem]{Remark}
\newtheorem{problem}[theorem]{Problem}

\newcommand{\R}{\mathbb{R}}
\newcommand{\E}{\mathbb{E}}
\newcommand{\PP}{\mathbb{P}}
\newcommand{\emp}{\bar\mu_N}
\newcommand{\vol}{\operatorname{vol}}

\title{The logarithm is necessary: exact critical scaling of the\\
uniform doubling condition on the ball}
\author{Yitzchak Shmalo\\
\small Einstein Institute of Mathematics, The Hebrew University of Jerusalem, Jerusalem, Israel\\
\small }
\date{June 2026}

\begin{document}
\maketitle

\begin{abstract}
In two companion papers we proved that $N$ independent uniform samples from a
convex body satisfy the uniform doubling condition of
[Shmalo, \emph{Ann.\ Math.\ Artif.\ Intell.} (2024)] with high probability, and
that for untruncated slabs over a body with an interior rolling ball the scaling
factor can be taken $\ell_N=O\bigl((\log N/N)^{2/(n+1)}\bigr)$, while single
gap configurations force $\ell_N\ge cN^{-2/(n+1)}$ with probability bounded
below.  Two questions remained: whether the failure probability can be driven
to $1$, and whether the logarithm is an artifact of the Vapnik--Chervonenkis
union bound.  We answer both for the unit ball, $n\ge2$: with probability
tending to one, among $\asymp(N/\log N)^{(n-1)/(n+1)}$ disjoint boundary caps
at least one contains a single, isolated, extremal sample point whose
surrounding slabs stall at empirical mass $1/N$ through a full grid step at
width scale $(\varepsilon\log N/N)^{2/(n+1)}$.  Consequently the uniform
doubling condition for pure slabs fails, with probability $1-o(1)$, for every
scaling factor $\ell\le c(\kappa)\,(\log N/N)^{2/(n+1)}$.  Combined with the
companion upper bound, the critical scaling of the uniform doubling condition
on the ball is
\[
\ell_N\;\asymp\;\Bigl(\frac{\log N}{N}\Bigr)^{2/(n+1)},
\]
exactly, logarithm included.  The proof is a second-moment argument over the
cap family; the only geometric input is the two-sided spherical cap-volume
asymptotic.
\end{abstract}

\section{Introduction}\label{sec:intro}

Let $K=B^n\subset\R^n$ be the closed unit ball, $n\ge2$, let
$\mu$ be the uniform probability measure on $K$, and let
$T=\{X_1,\dots,X_N\}$ be i.i.d.\ samples with empirical measure $\emp$.  For a
unit vector $u$, a center $s$, and a width $w$, write
$S_w=S_w(u,s)=\{x\in\R^n:|\langle x-s,u\rangle|\le w/2\}$ for the (pure,
untruncated) slab.  The uniform doubling condition (UDC) of
\cite{ShmaloUDC} with constants $(\kappa,\delta,\sigma,\beta,\ell)$ requires
\[
\emp(S_{\ell\kappa^{i+1}})\;\ge\;
\min\bigl\{\delta,\ (1+\sigma)\,\emp(S_{\ell\kappa^{i}})\bigr\}
\]
for every slab and every $i\in\mathbb N_0$ with $\ell\kappa^i<\beta$; it is the
hypothesis of the stability-of-accuracy theorems of \cite{BJS,ShmaloUDC}.  In
\cite{companion2} we showed that i.i.d.\ samples satisfy the UDC with
probability $1-\frac1N$ and scaling factor
$\ell_N=O((\log N/N)^{1/n})$, for all uniformly thick truncated-slab families
at once; in \cite{companion3} we showed that for \emph{pure} slabs over a body
with an interior rolling ball the scaling factor improves to
\[
\ell_N\;=\;O\Bigl(\Bigl(\frac{\log N}{N}\Bigr)^{2/(n+1)}\Bigr),
\]
that this exponent is dictated by the spherical cap volume, and that
truncations genuinely shift the critical exponent to $\frac1n$.  The
lower-bound constructions in both papers produce a \emph{single} violating
configuration, hence only a failure probability bounded below by a constant,
and at scale $N^{-2/(n+1)}$ --- without the logarithm.  This left
\cite[Problem 5.1]{companion3}:

\begin{quote}
\emph{Does the UDC for pure slabs on the ball hold at scaling
$\ell\asymp N^{-2/(n+1)}$, or is the logarithmic factor necessary?}
\end{quote}

We prove that the logarithm is necessary, and that the failure below the
logarithmic scale is overwhelming rather than merely non-negligible.

\begin{theorem}\label{thm:main}
Let $n\ge2$, let $\mu$ be uniform on the unit ball, and fix
$\kappa>1$, $\sigma>0$, $\beta>0$, and $\delta\ge\frac3N$.  There are
constants $\varepsilon_0=\varepsilon_0(n,\kappa)>0$ and
$c=c(n,\kappa)>0$ such that, with probability $1-o(1)$ as
$N\to\infty$, the doubling inequality fails for some pure slab at every grid
scale below $c\,(\log N/N)^{2/(n+1)}$.  Consequently, with probability
$1-o(1)$, the sample satisfies the UDC for pure slabs with \emph{no} scaling
factor
\[
\ell\;\le\;c\,\Bigl(\frac{\log N}{N}\Bigr)^{2/(n+1)} .
\]
Together with the upper bound of \cite{companion3}, the critical scaling of
the uniform doubling condition for pure slabs on the ball is
$\Theta\bigl((\log N/N)^{2/(n+1)}\bigr)$: both the exponent and the logarithm
are exact.
\end{theorem}

The interpretation is pleasingly sharp.  The scaling factor $\ell$ is the
resolution below which the doubling condition does not look; the theorem says
that for i.i.d.\ data this resolution cannot be chosen finer than the scale at
which the \emph{most isolated boundary point} of the sample sits alone --- and
on the ball, among the $\asymp(N/\log N)^{(n-1)/(n+1)}$ essentially
independent boundary caps, the most isolated one is isolated to depth
$\asymp(\log N/N)^{2/(n+1)}$, a factor $(\log N)^{2/(n+1)}$ deeper than a
typical extreme point.  The companion upper bound shows this is the only
obstruction.

\section{Preliminaries}\label{sec:prelim}

Throughout, $c_4(n)\le c_5(n)$ denote the spherical cap-volume constants: for
the cap $C(u,h):=\{x\in B^n:\langle x,u\rangle>1-h\}$,
\begin{equation}\label{eq:cap}
c_4\,h^{\frac{n+1}2}\;\le\;\Lambda(h):=\mu\bigl(C(u,h)\bigr)\;\le\;
c_5\,h^{\frac{n+1}2},
\qquad 0<h\le\tfrac12,
\end{equation}
by the standard computation
$\Lambda(h)=\frac{1}{\vol(B^n)}\int_{1-h}^1\omega_{n-1}(1-t^2)^{\frac{n-1}2}dt$
and $1-t^2\asymp(1-t)$ on $[\frac12,1]$.

We also record the violating-slab mechanism from
\cite[Theorem 1.3]{companion3}, specialized to pure slabs on the ball.

\begin{lemma}[Stalled slab]\label{lem:stall}
Fix $\kappa>1$, $\sigma>0$, $\delta\ge\frac3N$.  Let $u$ be a unit vector,
$0<h<H\le\frac12$ with $h\le\frac H{4}$, and suppose the sample satisfies:
exactly one point $X^*\in C(u,H)$, and moreover $X^*\in C(u,h)$.  Then for
every $w\le H/\kappa$, the slab $S_w(u,X^*)$ satisfies
\[
\emp\bigl(S_w(u,X^*)\bigr)=\emp\bigl(S_{\kappa w}(u,X^*)\bigr)=\frac1N
\;<\;\min\Bigl\{\delta,\ (1+\sigma)\,\emp\bigl(S_w(u,X^*)\bigr)\Bigr\},
\]
i.e.\ the doubling inequality fails at every scale $w\le H/\kappa$.
Consequently, if some grid scale $\ell\kappa^i$ ($i\in\mathbb N_0$,
$\ell\kappa^i<\beta$) lies in $(0,H/\kappa]$ --- which happens whenever
$\ell\le H/\kappa^2$ and $N$ is large --- the UDC with scaling factor $\ell$
fails.
\end{lemma}

\begin{proof}
Write $\tau(x)=\langle x,u\rangle$ and $\tau^*=\tau(X^*)\ge1-h$.  For
$w'\le H$,
\[
S_{w'}(u,X^*)\cap B^n\;\subseteq\;\{x\in B^n:\ \tau(x)\ge\tau^*-\tfrac{w'}2\}
\;\subseteq\;C\Bigl(u,\ h+\tfrac{w'}2\Bigr)\;\subseteq\;C(u,H),
\]
since $h+\frac{w'}2\le\frac H4+\frac H2<H$.  By hypothesis $C(u,H)$ contains
only the sample point $X^*$, which lies in every $S_{w'}(u,X^*)$; hence the
empirical masses of $S_w$ and $S_{\kappa w}$ both equal $\frac1N$ for
$w\le H/\kappa$.  Since $\delta\ge\frac3N>\frac1N$ and $\sigma>0$,
$\min\{\delta,(1+\sigma)\frac1N\}>\frac1N$, and the doubling inequality fails.
\end{proof}

\section{Proof of Theorem \ref{thm:main}}\label{sec:proof}

\begin{proof}[Proof of Theorem \ref{thm:main}]
Fix $\kappa,\sigma,\beta,\delta\ge\frac3N$ as in the statement, and set
\[
q:=\bigl(4\kappa^2\bigr)^{\frac{n+1}2},
\qquad
\varepsilon_0:=\frac{n-1}{2(n+1)}\cdot\frac{c_4}{c_5\,q},
\qquad
\tau_N:=\varepsilon_0\log N .
\]
Define the inner depth $h_N$ by $\Lambda(h_N)=\frac{\tau_N}{N}$ --- by
\eqref{eq:cap}, $h_N\asymp(\tau_N/N)^{2/(n+1)}$ --- and the outer depth
$H_N:=4\kappa^2h_N$, so that by \eqref{eq:cap},
\begin{equation}\label{eq:outer}
\Lambda(H_N)\;\le\;c_5(4\kappa^2)^{\frac{n+1}2}h_N^{\frac{n+1}2}
\;\le\;\frac{c_5\,q}{c_4}\cdot\frac{\tau_N}{N}
\;=:\;\frac{Q\tau_N}{N}.
\end{equation}

\emph{The cap family.}  The spherical cap $C(u,H_N)$ has angular radius
$\arccos(1-H_N)\le2\sqrt{H_N}$.  Choose a maximal family of unit vectors
$u_1,\dots,u_M$ with pairwise angles at least $8\sqrt{H_N}$; then the caps
$C(u_j,H_N)$ are pairwise disjoint, and a standard volumetric packing bound on
the sphere gives
\begin{equation}\label{eq:M}
M\;\ge\;c_6(n)\,H_N^{-\frac{n-1}2}
\;\ge\;c_7(n,\kappa)\Bigl(\frac{N}{\tau_N}\Bigr)^{\frac{n-1}{n+1}} .
\end{equation}

\emph{Per-cap events.}  Let
\[
E_j:=\Bigl\{\#\bigl(T\cap C(u_j,H_N)\bigr)=1\ \text{ and }\
T\cap C(u_j,H_N)\subseteq C(u_j,h_N)\Bigr\},
\]
the event that the $j$-th outer cap contains exactly one sample point and that
this point is $h_N$-deep.  Then
\begin{equation}\label{eq:p}
p:=\PP(E_j)=N\,\Lambda(h_N)\bigl(1-\Lambda(H_N)\bigr)^{N-1}
\;\ge\;\tau_N\,e^{-2Q\tau_N}
\;=\;\tau_N\,N^{-2Q\varepsilon_0},
\end{equation}
for large $N$, using $1-x\ge e^{-2x}$ for $0\le x\le\frac12$ and
\eqref{eq:outer}.

\emph{First moment.}  Let $X:=\sum_{j\le M}\mathbf 1_{E_j}$.  By
\eqref{eq:M} and \eqref{eq:p},
\[
\E X\;=\;Mp\;\ge\;c_7\Bigl(\frac{N}{\tau_N}\Bigr)^{\frac{n-1}{n+1}}
\tau_N\,N^{-2Q\varepsilon_0}
\;\ge\;c_7\,N^{\frac{n-1}{n+1}-2Q\varepsilon_0}\,
\tau_N^{\,1-\frac{n-1}{n+1}} .
\]
By the choice of $\varepsilon_0$,
$2Q\varepsilon_0=\frac{n-1}{n+1}\cdot\frac{c_4}{c_5}\le\frac{n-1}{n+1}$, with
strict inequality unless $c_4=c_5$; to be safe replace $\varepsilon_0$ by
$\varepsilon_0/2$, so that
$2Q\varepsilon_0\le\frac{n-1}{2(n+1)}$ and
\begin{equation}\label{eq:EX}
\E X\;\ge\;c_7\,N^{\frac{n-1}{2(n+1)}}\;\longrightarrow\;\infty .
\end{equation}

\emph{Second moment.}  For $i\ne j$ the caps are disjoint, and
\[
\PP(E_i\cap E_j)
=N(N-1)\,\Lambda(h_N)^2\bigl(1-2\Lambda(H_N)\bigr)^{N-2}.
\]
Hence
\[
\frac{\PP(E_i\cap E_j)}{p^2}
=\frac{N-1}{N}\cdot
\frac{\bigl(1-2\Lambda(H_N)\bigr)^{N-2}}{\bigl(1-\Lambda(H_N)\bigr)^{2N-2}}
\;\le\;\exp\Bigl(C\,N\,\Lambda(H_N)^2\Bigr)
\;=\;\exp\Bigl(O\Bigl(\frac{(\log N)^2}{N}\Bigr)\Bigr)
\;=\;1+o(1),
\]
uniformly in $i\ne j$, since
$\log\frac{(1-2x)}{(1-x)^2}=O(x^2)$ for $0\le x\le\frac14$.  Therefore
\[
\operatorname{Var}(X)\;=\;\E X+\sum_{i\ne j}\bigl(\PP(E_i\cap E_j)-p^2\bigr)
\;\le\;\E X+o(1)\,(\E X)^2,
\]
and by Chebyshev's inequality,
\[
\PP(X=0)\;\le\;\frac{\operatorname{Var}(X)}{(\E X)^2}
\;\le\;\frac1{\E X}+o(1)\;=\;o(1)
\]
by \eqref{eq:EX}.

\emph{Conclusion.}  With probability $1-o(1)$ some $E_j$ occurs.  On $E_j$,
Lemma~\ref{lem:stall} (with $u=u_j$, $h=h_N$, $H=H_N$; note
$h_N=H_N/(4\kappa^2)\le H_N/4$) shows the doubling inequality fails at every
scale $w\le H_N/\kappa$, so the UDC fails for every scaling factor
$\ell\le H_N/\kappa^2=4h_N$.  Since
$h_N\ge(c_5^{-1}\tau_N/N)^{2/(n+1)}\ge c(n,\kappa)(\log N/N)^{2/(n+1)}$, the
theorem follows.
\end{proof}

\begin{remark}[General rolling-ball bodies]
The proof used the ball only through the two-sided cap asymptotic
\eqref{eq:cap} and the packing bound \eqref{eq:M}.  Both hold, with constants
depending on the curvature bounds, for any convex body with $C^2$ boundary and
principal curvatures bounded between two positive constants: caps satisfy
$\Lambda(h)\asymp h^{(n+1)/2}$ uniformly in the direction, and a maximal
separated family of boundary directions provides
$\asymp H^{-(n-1)/2}$ disjoint caps.  Theorem~\ref{thm:main} therefore extends
verbatim to this class, and with the upper bound of \cite{companion3} the
critical scaling $\Theta((\log N/N)^{2/(n+1)})$ holds for all such bodies.
Densities bounded above and below change only the constants.
\end{remark}

\begin{remark}[The truncated case]
The multiplication-of-configurations mechanism should also upgrade the
truncated lower bound of \cite{companion3} from constant probability at scale
$N^{-1/n}$ to probability $1-o(1)$ at scale $(\log N/N)^{1/n}$, matching the
upper bound of \cite{companion2} exactly: since the UDC is quantified over all
truncation sets of a given size, one may place
$\asymp h_N^{-(n-1)/2}$ rotated copies of a thick conical truncation region
with apices along the boundary of a smooth body and run the same second-moment
argument over the resulting cone caps, whose masses scale as $h^{n}$.  The
geometric bookkeeping (uniform thickness of the rotated regions and
disjointness of the cone caps) is routine but not carried out here.
\end{remark}

\begin{problem}
Determine the exact constant: does
$\ell_N^{\mathrm{crit}}/(\log N/N)^{2/(n+1)}$ converge, and to what, as
$N\to\infty$?  The upper route loses a factor through the VC bound, the lower
route through the choice $H=4\kappa^2h$; we expect a clean extremal constant
expressible through the cap-volume constant $c_4=c_4(n)$ and $\kappa$ alone.
\end{problem}

\subsection*{Acknowledgments}
The research presented in this paper was supported by the European Research
Council (ERC) under the European Union's Horizon Europe research and innovation
programme (grant agreement No.\ 101041711), by the Israel Science Foundation
(grant number 2258/19) and by the Israel Science Foundation (ISF Grant
4101/25).

\subsection*{Declaration on generative AI}
During the preparation of this work the author used generative AI tools to
accelerate drafting and revision.  All mathematical claims, proofs, and
bibliographic information were subsequently reviewed and edited by the author,
who takes full responsibility for the content of the manuscript.

\begin{thebibliography}{9}

\bibitem{BJS}
L.~Berlyand, P.-E.~Jabin and C.~A.~Safsten,
\emph{Stability for the training of deep neural networks and other
classifiers},
Math.\ Models Methods Appl.\ Sci.\ \textbf{31} (2021), no.~11, 2345--2390.

\bibitem{ShmaloUDC}
Y.~Shmalo,
\emph{Stability of accuracy for the training of DNNs via the uniform doubling
condition},
Ann.\ Math.\ Artif.\ Intell.\ (2024), DOI 10.1007/s10472-023-09919-1.

\bibitem{companion2}
Y.~Shmalo,
\emph{The uniform doubling condition holds with high probability for
independent samples},
manuscript, 2026.

\bibitem{companion3}
Y.~Shmalo,
\emph{The critical scaling of the uniform doubling condition: smooth bodies
versus truncations},
manuscript, 2026.

\end{thebibliography}

\end{document}
